\documentclass[11pt]{amsart}
\usepackage[margin=1.1in]{geometry}
\usepackage{amsmath,amssymb,amsthm,mathtools}
\usepackage{booktabs,longtable}
\usepackage{enumitem}
\usepackage[hidelinks]{hyperref}

\newtheorem{theorem}{Theorem}[section]
\newtheorem{lemma}[theorem]{Lemma}
\newtheorem{proposition}[theorem]{Proposition}
\newtheorem{corollary}[theorem]{Corollary}
\newtheorem{question}[theorem]{Question}
\theoremstyle{definition}
\newtheorem{definition}[theorem]{Definition}
\theoremstyle{remark}
\newtheorem{remark}[theorem]{Remark}

\newcommand{\Pl}{P_\ell}
\newcommand{\1}{\mathbf 1}

\title{The triangular prism is enumeratively chromatic-choosable}
\author{Dylan Tague}
\address{Thomas Edison State University, Trenton, New Jersey, USA}
\email{dylan.tague@students.tesu.edu}
\subjclass[2020]{05C15 (primary); 05C30, 05C31 (secondary)}
\keywords{list colouring, list colour function, chromatic polynomial, enumerative chromatic-choosability,
line graph, Galvin's theorem, Cartesian product}
\date{September 27, 2026}

\begin{document}

\begin{abstract}
For a graph $G$ and a positive integer $m$, let $P(G,m)$ be the chromatic polynomial and let
$\Pl(G,m)$ be the minimum number of proper colourings from an assignment of $m$-element lists.
A graph is \emph{enumeratively chromatic-choosable} if $\Pl(G,m)=P(G,m)$ for every $m$.
We prove that the triangular prism $K_3\square K_2$ has this property. Since
$K_3\square K_2$ is the line graph of $K_{2,3}$, the statement is equivalent to an
enumerative strengthening of Galvin's theorem for $K_{2,3}$: for every $k\ge 3$ and every
assignment of $k$-element lists to the six edges of $K_{2,3}$, the number of proper list
edge-colourings is at least $k(k-1)(k-2)(k^3-6k^2+14k-13)$, the number of proper
$k$-edge-colourings. General threshold results of Dong and Zhang already give the equality for
$m\ge 8$, so it remains here to establish $3\le m\le 7$. The result is a positive instance, for the
pair $(K_3,K_2)$, of a question of Kirov and Naimi on Cartesian products of $n$-monophilic graphs.
The proof reduces the problem to a sharp counting inequality for list colourings of a five-vertex
``house'' graph, which has a unique extremal list configuration. Two essential steps are computer-assisted
and verified by exact arithmetic: linear-programming dual certificates, and a finite enumeration of
membership types for small list sizes. Routine polynomial identities elsewhere were also verified
symbolically.
\end{abstract}

\maketitle

% =====================================================================
\section{Introduction}\label{sec:intro}

Let $G$ be a finite simple graph. A \emph{list assignment} $L$ gives each vertex $v$ a finite set
$L(v)$ of colours, and an $L$-colouring is a proper colouring $\varphi$ with $\varphi(v)\in L(v)$
for all $v$. Write $P(G,L)$ for the number of $L$-colourings. If every list has size $m$, then $L$
is an \emph{$m$-assignment}. The \emph{list colour function} is
\[
\Pl(G,m)=\min\{P(G,L): L \text{ an $m$-assignment of } G\}.
\]
Giving every vertex the same list shows that $\Pl(G,m)\le P(G,m)$, where $P(G,m)$ is the chromatic
polynomial. The function $\Pl$ was introduced by Kostochka and Sidorenko~\cite{KS}. Donner~\cite{Donner}
proved that $\Pl(G,m)=P(G,m)$ for all sufficiently large $m$. The threshold for this equality has been
bounded by Thomassen~\cite{Thomassen}, by Wang, Qian and Yan~\cite{WQY}, by Dong and
Zhang~\cite{DZ23}, and, linearly in the maximum degree, by Zhang and Dong~\cite{ZD26}.

Following Allred and Mudrock~\cite{AM} and Chi et al.~\cite{CLMNW}, a graph $G$ is
\emph{enumeratively chromatic-choosable} if $\Pl(G,m)=P(G,m)$ for every $m\in\mathbb N$. Equality at
one fixed $m$ is the older notion of an \emph{$m$-monophilic} graph, due to Kirov and Naimi~\cite{KN}.
Zhang and Dong~\cite{ZD26np} showed that equality at $m$ need not persist to $m+1$, so equality has to
be established at every $m$ separately. Allred and Mudrock~\cite[Theorem~4]{AM} showed that chordal
graphs, cycles and $K_{2,3}$ are enumeratively chromatic-choosable, and characterised the enumeratively
chromatic-choosable graphs of chromatic number two. Chi et al.~\cite{CLMNW} studied theta graphs. All of
these results concern vertex colourings of the graphs named; in particular, the result for $K_{2,3}$
concerns $K_{2,3}$ itself and not its line graph, which is the prism studied here.

\subsection*{Line graphs and a question of Kirov and Naimi}
Galvin~\cite{Galvin} proved that every bipartite multigraph $R$ has list chromatic index $\Delta(R)$. In
other words, the line graph $L(R)$ is $\chi(L(R))$-choosable. This settled Dinitz's problem, which is
the case $R=K_{n,n}$. Kirov and Naimi observed that $L(K_{n,n})$ is the Cartesian product
$K_n\square K_n$, and proposed the following ``other way to generalize the Dinitz Conjecture''.

\begin{question}[{Kirov--Naimi \cite[Question~1]{KN}}]\label{q:KN}
Is the product of two $n$-monophilic graphs $n$-monophilic?
\end{question}

Here, as in the preceding observation, the product is the Cartesian product.

They showed that the answer is negative for $n=2$, the product $P_2\square P_3$ being a counterexample,
and raised the possibility that it is positive for $n\ge 3$. As far as we know the question is open in
general. A machine-checked formalization project of Kirov~\cite{KirovLean} records it as open and proves
that $\Pl(G,m)=P(G,m)$ for every $m\ge3$ when $G$ is a ladder $P_2\square P_n$ or a grid $P_3\square P_n$.
These give positive instances in which both factors are paths. Among products of two nontrivial
complete graphs, $K_2\square K_2=C_4$ is a cycle and is already covered. A natural first non-cycle case is
the \emph{triangular prism} $K_3\square K_2$, which is also the line graph $L(K_{2,3})$.

\subsection*{Main result}
\begin{theorem}\label{thm:main}
The triangular prism $K_3\square K_2$ is enumeratively chromatic-choosable. Equivalently, for every
$k\ge 3$ and every assignment $L$ of $k$-element lists to the edges of $K_{2,3}$, the number
$N_L(K_{2,3})$ of proper $L$-edge-colourings satisfies
\begin{equation}\label{eq:main}
N_L(K_{2,3})\ \ge\ c_k:=k(k-1)(k-2)\,(k^3-6k^2+14k-13),
\end{equation}
and $c_k$ is the number of proper $k$-edge-colourings of $K_{2,3}$.
\end{theorem}

In a companion paper~\cite{TagueExact} the author proved that the number of list edge-colourings of a
left-$k$-regular bipartite multigraph $R$ is at least $c_k(R)$, the number of proper
$k$-edge-colourings, in the special case of \emph{exact left lists}: all edges at a vertex of one shore
share a list of size $k$, and every entry must be used. That paper asks~\cite[Question~11.4]{TagueExact} whether the inequality survives for arbitrary per-edge lists.
Theorem~\ref{thm:main} answers this question for $R=K_{2,3}$. Its proof is independent of the methods
of~\cite{TagueExact} and uses no result from it.

For $m<3$ both $P(K_3\square K_2,m)$ and $\Pl(K_3\square K_2,m)$ are zero, because the prism contains
a triangle; so the two formulations agree. Since $K_3$ and $K_2$ are chordal, and hence $3$-monophilic,
Theorem~\ref{thm:main} answers Question~\ref{q:KN} positively for the particular pair
$(K_3,K_2)$. It gives no information about the question for other pairs.

\begin{remark}[What was known]\label{rem:known}
The prism has $9$ edges. Dong and Zhang~\cite{DZ23} proved $\Pl(G,m)=P(G,m)$ for
$m\ge |E(G)|-1$, which covers the prism for every $m\ge 8$. Kirov's formalization~\cite{KirovLean}
proves equality for every subcubic graph when $m\ge 15$, and Zhang and Dong~\cite{ZD26} give
$m\ge 23.41\,\Delta(G)$. Theorem~\ref{thm:main} therefore covers the remaining list sizes $m=3,\dots,7$, including
the critical value $m=\chi=3$. In our searches we found no earlier source proving any of these cases;
we have not certified novelty beyond those searches.
\end{remark}

\subsection*{Strategy}
Fix the colour $x$ of one edge $uw$ incident with a degree-$3$ vertex. Summing over $x$ writes
$N_L(K_{2,3})$ as a sum of row counts $S(x)$ (Section~\ref{sec:reduction}). The \emph{deficit} of a row
is the amount by which it falls short of its common-palette value. The published fact that cycles are
enumeratively chromatic-choosable disposes of $k=3$ and of every row whose colour is absent from the list
of the other edge at $w$. Each remaining row is a list-colouring count of a five-vertex ``house''
graph. The core of the paper is the \emph{House Lemma}
(Theorem~\ref{thm:house}). It bounds the house count from below by its common-palette value minus
$2(k-2)(k-3)$, and it shows that falling below the common value forces one explicit list
configuration $\Delta$. We then compute the deficits of the other rows in configuration $\Delta$
exactly and show that they pay for the unique possible deficit (Section~\ref{sec:pair}). Summing
completes the proof (Section~\ref{sec:proof}).

\subsection*{Computer assistance}
The reduction, the case of equal list pairs (Proposition~\ref{prop:equalpairs}), the localization and
compensation arguments, and the final assembly are ordinary proofs. Two essential steps use a
computer, in precisely identified ways:
\begin{enumerate}[label=(\alph*)]
\item for unequal shoulder lists and $k\ge 6$ (Proposition~\ref{prop:ues}), explicit
  linear-programming dual certificates, whose $71$ defining inequalities are verified exactly. For
  $k\ge9$ the certificate is given by four cubic polynomials, and for $k=6,7,8$ by integer vectors;
\item for $k\in\{4,5\}$, a complete enumeration of membership types, which covers every list
  configuration in every colour universe (Section~\ref{sec:computation}). It has $49{,}394$ and
  $519{,}800$ cases.
\end{enumerate}
Routine polynomial identities elsewhere were also verified symbolically. The most substantial of these is
the expansion of an explicit quadratic polynomial in the equal-shoulder case
(Proposition~\ref{prop:es}), whose linear coefficients are printed in Table~\ref{tab:esLinear}. All
computations are exact and have been reimplemented independently; see
Section~\ref{sec:computation}.

% =====================================================================
\section{Preliminaries}\label{sec:prelim}

\subsection{The graph and its common count}
Let $K_{2,3}$ have degree-$3$ vertices $u,v$ and degree-$2$ vertices $a,b,w$. Name its edges
\[
e_1=ua,\quad e_2=av,\quad e_3=ub,\quad e_4=bv,\quad e_5=uw,\quad e_6=wv .
\]
In the line graph, $\{e_1,e_3,e_5\}$ and $\{e_2,e_4,e_6\}$ are triangles, and the three ``rungs''
$e_1e_2$, $e_3e_4$, $e_5e_6$ join them. This is the prism $K_3\square K_2$. Throughout, an edge
colouring means a proper edge colouring, and lists are attached to edges.

\begin{lemma}\label{lem:ck}
The number of proper $k$-edge-colourings of $K_{2,3}$ is $c_k=(k)_3\,(k^3-6k^2+14k-13)$, where
$(k)_3=k(k-1)(k-2)$.
\end{lemma}
\begin{proof}
Colour the triangle $e_1,e_3,e_5$ injectively, in $(k)_3$ ways, say with colours $\gamma_1,\gamma_3,\gamma_5$.
The triangle $e_2,e_4,e_6$ must then be coloured injectively with $e_2\ne\gamma_1$, $e_4\ne\gamma_3$ and
$e_6\ne\gamma_5$. By inclusion--exclusion over these three conditions the number of ways is
$(k)_3-3(k-1)(k-2)+3(k-2)-1=k^3-6k^2+14k-13$.
\end{proof}

We write
\begin{gather*}
o_k:=\frac{c_k}{k(k-1)}=(k-2)(k^3-6k^2+14k-13),\\
d_k:=(k-1)(k-2)(k^2-5k+7),\qquad \sigma_k:=2(k-2)(k-3).
\end{gather*}
Note that $d_k=(k-2)^4+(k-2)=P(C_4,k-1)$, that $o_k-d_k=\sigma_k$, and that $d_3=o_3=2$.

\subsection{Two basic facts}
\begin{enumerate}[label=(F\arabic*)]
\item\label{F:mono} \emph{Monotonicity.} Enlarging lists cannot decrease the number of list colourings.
In particular, if $|L(v)|\ge m$ for every vertex, then $P(G,L)\ge \Pl(G,m)$.
\item\label{F:cycle} \emph{Cycles.} Cycles are enumeratively chromatic-choosable
\cite[Theorem~4(ii)]{AM}. In particular, if $|L(v)|\ge m$ for every vertex of $C_4$, then
$P(C_4,L)\ge (m-1)^4+(m-1)$.
\end{enumerate}

% =====================================================================
\section{Row kernels and deficits}\label{sec:reduction}

Fix a $k$-assignment $L$ of $E(K_{2,3})$, with $k\ge3$, and put $A=L(e_5)$ and $B=L(e_6)$. The edges
$e_1,e_2,e_3,e_4$ form a $4$-cycle $Q$ through $u$ and $v$. For colours $x,y$, let $F(x,y)$ be the
number of proper colourings of $Q$ from its lists in which $x$ appears on neither $e_1$ nor $e_3$ (the
edges of $Q$ at $u$), and $y$ appears on neither $e_2$ nor $e_4$ (the edges of $Q$ at $v$).

\begin{lemma}\label{lem:glue}
$N_L(K_{2,3})=\sum_{x\in A}S(x)$, where $S(x):=\sum_{y\in B\setminus\{x\}}F(x,y)$.
\end{lemma}
\begin{proof}
A colouring of $K_{2,3}$ is determined by the colours $x$ of $e_5$ and $y$ of $e_6$, with $x\neq y$
because the two edges share $w$, together with a colouring of $Q$. The only other conflicts are between
$e_5$ and the $Q$-edges at $u$, and between $e_6$ and the $Q$-edges at $v$.
\end{proof}

For the common palette, $F(x,y)$ takes the same value $F^*$ for all $x\ne y$, by symmetry. Applying
Lemma~\ref{lem:glue} to the common assignment gives $k(k-1)F^*=c_k$, so $F^*=o_k$. For the given
assignment $L$ define the \emph{deficit} of a colour $x$ by
\[
\delta_x:=(k-1)o_k-S(x).
\]
(The quantity $\delta_x$ is defined for every colour $x$, but only the rows
$x\in A$ enter the count.) Then $N_L(K_{2,3})=c_k-\sum_{x\in A}\delta_x$, so Theorem~\ref{thm:main} is equivalent to
\begin{equation}\label{eq:deficits}
\sum_{x\in A}\delta_x\le 0 .
\end{equation}

\begin{lemma}[Cycle bound]\label{lem:cyclebound}
$F(x,y)\ge d_k$ for all colours $x,y$.
\end{lemma}
\begin{proof}
Delete $x$ from the lists of $e_1$ and $e_3$, and $y$ from the lists of $e_2$ and $e_4$. Each edge of
$Q$ is incident with exactly one of $u,v$, so each list loses at most one colour and keeps at least
$k-1$. The line graph of $Q$ is $C_4$, so \ref{F:cycle} gives at least $P(C_4,k-1)=d_k$ colourings.
\end{proof}

\begin{corollary}\label{cor:easy}
\begin{enumerate}[label=(\roman*)]
\item If $k=3$, then $\delta_x\le 0$ for every $x\in A$. Hence Theorem~\ref{thm:main} holds for $k=3$.
\item If $x\in A\setminus B$, then $\delta_x\le -(k-1)(k-2)(k^2-7k+13)<0$.
\end{enumerate}
\end{corollary}
\begin{proof}
The sum defining $S(x)$ has at least $k-1$ terms, each at least $d_k$ by Lemma~\ref{lem:cyclebound}.
For $k=3$, $d_3=o_3$, which gives (i). If $x\notin B$, the sum has $k$ terms, so
$S(x)\ge k\,d_k$, and $k\,d_k-(k-1)o_k=(k-1)(k-2)(k^2-7k+13)$. The quadratic $k^2-7k+13$ has negative
discriminant, which gives (ii).
\end{proof}

\subsection{\texorpdfstring{Rows in $B$ are house counts}{Rows in B are house counts}}
The \emph{house} $H$ is the line graph of a $4$-cycle with a pendant edge. Its vertices are two
\emph{base} vertices $\beta_1,\beta_3$, two \emph{shoulder} vertices $s_2,s_4$ and a \emph{roof} $\rho$,
and its edges are $\beta_1\beta_3$, $\beta_1s_2$, $\beta_3s_4$, $s_2s_4$, $s_2\rho$ and $s_4\rho$.
A \emph{house assignment} is a tuple $(A_1,A_3,L_2,L_4,R)$ of lists on
$(\beta_1,\beta_3,s_2,s_4,\rho)$. We write $h(A_1,A_3,L_2,L_4,R)$ for its number of colourings.

\begin{lemma}\label{lem:rowhouse}
For $x\in A\cap B$, $S(x)=h\bigl(L(e_1)\setminus\{x\},\,L(e_3)\setminus\{x\},\,L(e_2),\,L(e_4),\,B\setminus\{x\}\bigr)$.
\end{lemma}
\begin{proof}
Summing over $y\in B\setminus\{x\}$ amounts to colouring a pendant edge at $v$, adjacent to $e_2$ and
$e_4$, from the list $B\setminus\{x\}$. Excluding $x$ at $u$ amounts to deleting $x$ from the lists of
$e_1$ and $e_3$. The resulting graph is the $4$-cycle $Q$ with a pendant edge at $v$, whose line graph
is $H$.
\end{proof}

In Lemma~\ref{lem:rowhouse} the roof list has size $k-1$, the shoulder lists have size $k$, and each
base list has size $k-1$ or $k$.

% =====================================================================
\section{The House Lemma}\label{sec:house}

\begin{theorem}[House Lemma]\label{thm:house}
Let $k\ge4$, and let $(A_1,A_3,L_2,L_4,R)$ be a house assignment with $|L_2|=|L_4|=k$, $|R|=k-1$ and
$|A_1|,|A_3|\in\{k-1,k\}$. Then
\[
h(A_1,A_3,L_2,L_4,R)\ \ge\ (k-1)o_k-\sigma_k .
\]
If moreover $h<(k-1)o_k$, then $|A_1|=|A_3|=k-1$ and the assignment is the configuration $\Delta$:
there are a $(k-1)$-set $C$, a colour $\gamma\notin C$ and a colour $c_0\in C$ with
\[
A_1=A_3=C,\qquad L_2=L_4=C\cup\{\gamma\},\qquad R=(C\setminus\{c_0\})\cup\{\gamma\}.
\]
\end{theorem}

The proof occupies the rest of this section. Write $\mathrm{target}=(k-1)o_k$.

\subsection{Reduction to exact sizes}
\begin{lemma}\label{lem:size}
If $|A_1|=k$ or $|A_3|=k$, then $h\ge\mathrm{target}$, provided the House Lemma holds for base lists
of size exactly $k-1$.
\end{lemma}
\begin{proof}
Suppose $|A_1|=k$. Choose any $(k-1)$-subset $C_3\subseteq A_3$. Since $|A_1|>|C_3|$, there is a colour
$a\in A_1\setminus C_3$. Let $C_1$ consist of $a$ and any $k-2$ further elements of $A_1$. Then
$|C_1|=|C_3|=k-1$ and $C_1\ne C_3$ because $a\in C_1\setminus C_3$. By \ref{F:mono},
$h(A_1,A_3,L_2,L_4,R)\ge h(C_1,C_3,L_2,L_4,R)$. The right-hand assignment has unequal base lists, so it
is not $\Delta$, and by the exact-size case its count is at least the target. The case $|A_3|=k$ is
symmetric.
\end{proof}

By Lemma~\ref{lem:size} it suffices to prove Theorem~\ref{thm:house} when
\begin{equation}\label{eq:sizes}
|A_1|=|A_3|=|R|=k-1,\qquad |L_2|=|L_4|=k,
\end{equation}
which we assume from now on. Put $I=A_1\cap A_3$ and write $\1_X$ for the indicator function of a set $X$.

\subsection{The shoulder-first formula}
\begin{lemma}\label{lem:shoulder}
Under \eqref{eq:sizes},
\begin{multline}\label{eq:shoulder}
h=\sum_{\substack{c_2\in L_2,\ c_4\in L_4\\ c_2\neq c_4}}
\bigl(k-1-\1_R(c_2)-\1_R(c_4)\bigr)\\
\times\Bigl[(k-1-\1_{A_1}(c_2))(k-1-\1_{A_3}(c_4))-|I|+\1_I(c_2)+\1_I(c_4)\Bigr].
\end{multline}
\end{lemma}
\begin{proof}
Colour the shoulders first, with $c_2\ne c_4$. The roof must avoid both, which leaves
$k-1-\1_R(c_2)-\1_R(c_4)$ choices. The base vertex $\beta_1$ must avoid $c_2$, and $\beta_3$ must avoid
$c_4$ and the colour of $\beta_1$. The ordered pairs from $(A_1\setminus c_2)\times(A_3\setminus c_4)$
number $(k-1-\1_{A_1}(c_2))(k-1-\1_{A_3}(c_4))$. Among them, those with equal entries correspond to the
colours of $I\setminus\{c_2,c_4\}$, and there are $|I|-\1_I(c_2)-\1_I(c_4)$ of these, because
$c_2\neq c_4$.
\end{proof}

Only the colours of $U=L_2\cup L_4$ enter \eqref{eq:shoulder} individually. Colours outside $U$ enter
only through $|I|$. Moreover, $h$ is affine in $|I|$ with coefficient
$-\sum_{c_2\ne c_4}(k-1-\1_R(c_2)-\1_R(c_4))\le -(k^2-k)(k-3)<0$. We therefore have the following.

\begin{lemma}\label{lem:Ibound}
Let $h^{\sharp}$ be the value of \eqref{eq:shoulder} with $|I|$ replaced by
$|I\cap U|+\min\{|A_1\setminus U|,|A_3\setminus U|\}$. Then $h\ge h^{\sharp}$.
\end{lemma}
\begin{proof}
We have $I\setminus U\subseteq A_1\setminus U$ and $I\setminus U\subseteq A_3\setminus U$, so $|I|$ is at
most the substituted value. The coefficient of $|I|$ is negative.
\end{proof}

We treat three cases: equal list pairs, equal shoulders with unequal bases, and unequal shoulders.

\subsection{\texorpdfstring{Equal pairs: $L_2=L_4$ and $A_1=A_3$}{Equal pairs}}
\begin{proposition}\label{prop:equalpairs}
Let $L_2=L_4=V$ and $A_1=A_3=A$ satisfy \eqref{eq:sizes}, and put $r=|A\cap V|\in\{0,\dots,k-1\}$. Then
$h\ge\mathrm{target}-\sigma_k$. Moreover $h<\mathrm{target}$ only when $r=k-1$ and $R$ consists of
the unique colour $\gamma$ of $V\setminus A$ together with $k-2$ colours of $A$, that is, only in
configuration $\Delta$; there $h=\mathrm{target}-\sigma_k$.
\end{proposition}
\begin{proof}
Condition on the roof colour $y\in R$. The remaining $4$-cycle has base lists $A$ and shoulder lists
$V\setminus\{y\}$. Its colourings are ordered pairs of distinct base colours and ordered pairs of distinct
shoulder colours, excluding coincidences across the two rungs. By inclusion--exclusion over those two
coincidences, $y$ contributes
\[
f(y)=p(p-1)q(q-1)-2t(p-1)(q-1)+t(t-1),\qquad p=k-1,\ q=|V\setminus\{y\}|,\ t=|(A\cap V)\setminus\{y\}|,
\]
and $h=\sum_{y\in R}f(y)$. There are three types of colour $y$:
\[
\begin{aligned}
f_{C}&=(k-1)^2(k-2)^2-2(r-1)(k-2)^2+(r-1)(r-2) &&(y\in A\cap V),\\
f_{V}&=(k-1)^2(k-2)^2-2r(k-2)^2+r(r-1) &&(y\in V\setminus A),\\
f_{O}&=k(k-1)^2(k-2)-2r(k-1)(k-2)+r(r-1) &&(y\notin V).
\end{aligned}
\]
Here $f_C-f_V=2\bigl((k-2)^2-(r-1)\bigr)\ge0$. Also $f_O-f_C$ is affine in $r$ and positive at $r=0$ and at
$r=k-1$ for $k\ge4$, so $f_O\ge f_C\ge f_V$ for all $r$. Hence $h$ is minimised by choosing $R$ greedily: first the $k-r$ colours of
$V\setminus A$, then colours of $A\cap V$. For $r\ge1$ this gives $h\ge(k-r)f_V+(r-1)f_C$, and for $r=0$
it gives $h\ge(k-1)f_V$. The maximal deficit $D(r)=\mathrm{target}-\bigl((k-r)f_V+(r-1)f_C\bigr)$
satisfies:
\begin{itemize}
\item it is a quadratic in $r$ with leading coefficient $3-k<0$;
\item $D'(k-1)=(k-3)(2k^2-8k+9)>0$ for $k\ge4$, so $D$ is increasing on $[1,k-1]$;
\item $D(k-1)=\sigma_k$, while $D(k-2)=-2(k-3)(k^2-5k+7)<0$;
\item the $r=0$ value is $-(k-1)(k-2)(2k^2-9k+11)<0$.
\end{itemize}
Hence a deficit requires $r=k-1$, i.e.\ $A\subset V$ and $V=A\cup\{\gamma\}$. In that case the greedy
roof is $\{\gamma\}\cup(k-2 \text{ colours of } A)$. Any other roof either omits $\gamma$, which raises
$h$ by at least $f_C-f_V=\sigma_k$, or uses a colour outside $V$, which raises it by at least
$f_O-f_C=2(k-2)(k^2-4k+5)>\sigma_k$ (values at $r=k-1$). Either way the deficit is removed. All the displayed polynomial identities are
routine to check; see Section~\ref{sec:computation}.
\end{proof}

\subsection{Equal shoulders, unequal bases}
\begin{proposition}\label{prop:es}
If $L_2=L_4=V$ and $A_1\ne A_3$ satisfy \eqref{eq:sizes}, then
$h\ \ge\ \mathrm{target}+q_k$, where $q_k=k^3-8k^2+24k-25>0$.
Equality is attained when $A_1,A_3,R$ are $V$ minus three distinct colours.
\end{proposition}

\begin{proof}
Classify the $k$ colours of $V$ by their membership in $(R,A_1,A_3)$, writing $n_{\tau}$ for the number
of colours of type $\tau\in\{0,1\}^3$. We have $n_{111}=k-\sum_{\tau\neq111}n_\tau$. Every membership statistic entering
\eqref{eq:shoulder}, other than $|I|$, is linear in these counts. Swapping $A_1$ and $A_3$ is a symmetry
of $H$ when $L_2=L_4$, so we may assume $|A_1\cap V|\ge|A_3\cap V|$. By Lemma~\ref{lem:Ibound} we may
take $|I|=|I\cap V|+(k-1-|A_1\cap V|)$.

\emph{Case A: $A_1\cap V\ne A_3\cap V$.}
Substituting $k=4+m$, the quantity $F=h^\sharp-\mathrm{target}-q_k$ becomes a quadratic polynomial in
the seven counts $n_\tau$ ($\tau\ne111$), with coefficients in $\mathbb Z[m]$. Its constant term is
$-(5m^3+28m^2+54m+37)$, and \emph{every other coefficient is a polynomial in $m$ with nonnegative
coefficients}. Table~\ref{tab:esLinear} lists the linear coefficients. Each of the $25$ nonzero quadratic
coefficients has the form $am+b$ with integers $a\ge1$ and $b\ge0$. Hence $F$ is nondecreasing in each count. The constraints force three
conditions:
\begin{itemize}
\item an $A_1$-only colour in $V$ (a type $010$ or $110$), since $|A_1\cap V|\ge|A_3\cap V|$ and the
  two sets differ;
\item a colour of $V$ missing from $A_1$, since $|A_1|=k-1$;
\item a colour of $V$ missing from $R$, since $|R|=k-1$.
\end{itemize}
The minimal count vectors satisfying these hitting conditions are $0/1$-vectors supported on one of
eight sets. At each of them $F$ is a polynomial with nonnegative coefficients in $m$:
\[
\begin{array}{ll@{\qquad}ll}
\{011,101,110\}: & 0, & \{000,010\}: & 3k^3-15k^2+25k-12,\\
\{011,100,110\}: & (k-2)^3, & \{000,110\}: & k^3-4k^2+2k+5,\\
\{001,110\}: & (k-3)^2, & \{001,010\}: & 2(k^3-5k^2+8k-3),\\
\{010,101\}: & (k-3)^2, & \{010,100\}: & k^3-5k^2+6k+1.
\end{array}
\]

\emph{Case B: $A_1\cap V=A_3\cap V$ but $A_1\ne A_3$.}
Only the types $000,100,011,111$ occur, and $j\ge1$ colours of $A_1\setminus V$ are not in $A_3$. Then
$G=h-\mathrm{target}-q_k$ is linear in $j$, with slope $(k-1)(k^2-3k+2n_{000}+2n_{011})>0$. At $j=1$ it is
nondecreasing in $(n_{000},n_{100},n_{011})$, subject to $n_{000}+n_{100}\ge2$ and $n_{000}+n_{011}\ge1$. The first constraint holds because
$|A_1\setminus V|=n_{000}+n_{100}-1$ must be at least $j=1$. The second holds because $R$ misses a colour
of $V$.
The three minimal vectors give
\begin{gather*}
(1,1,0):\ 2k^3-10k^2+15k-5,\qquad (2,0,0):\ (k-1)(4k^2-16k+19),\\ (0,2,1):\ 2k^3-12k^2+25k-19,
\end{gather*}
all positive for $k\ge4$.
\end{proof}

\begin{table}[ht]
\centering
\small
\begin{tabular}{@{}ll@{\qquad}ll@{}}
\toprule
count & coefficient in $F$ ($k=4+m$) & count & coefficient in $F$ ($k=4+m$)\\
\midrule
$n_{110}$ & $2(m+1)(m+2)(m+3)$ & $n_{010}$ & $4m^3+25m^2+48m+30$\\
$n_{101}$ & $m(m+1)(m+3)$ & $n_{001}$ & $3m^3+17m^2+29m+18$\\
$n_{100}$ & $(m+1)(2m^2+8m+5)$ & $n_{000}$ & $4m^3+24m^2+43m+25$\\
$n_{011}$ & $2(m+3)(m^2+3m+3)$ & constant & $-(5m^3+28m^2+54m+37)$\\
\bottomrule
\end{tabular}
\caption{Linear coefficients of $F$ in Case A of Proposition~\ref{prop:es}.}
\label{tab:esLinear}
\end{table}

\subsection{Unequal shoulders}
\begin{proposition}\label{prop:ues}
If $L_2\ne L_4$ and \eqref{eq:sizes} holds, then $h\ge\mathrm{target}$ for every $k\ge 6$.
\end{proposition}

\begin{proof}[Proof outline]
Put $D=L_2\cap L_4$, $P=L_2\setminus L_4$ and $Q=L_4\setminus L_2$. Then $|P|=|Q|=b\ge1$, and we fix a
bijection $P\to Q$. The variables are:
\begin{itemize}
\item seven counts $D_\tau$ of colours of $D$ by type $\tau\in\{0,1\}^3\setminus\{111\}$, where types
  record membership in $(R,A_1,A_3)$;
\item $64$ counts $Z_{\tau,\tau'}$ of paired colours whose $P$-member has type $\tau$ and whose
  $Q$-member has type $\tau'$.
\end{itemize}
Then $D_{111}=k-\sum D_\tau-\sum Z$, so every variable lies in $[0,k]$, and $|L_4|=k$ holds automatically.
All intersection statistics in \eqref{eq:shoulder} are linear in these $71$ variables. Since
$|I\setminus U|\le|A_1\setminus U|=k-1-|A_1\cap U|$ and the coefficient of $|I|$ is negative, we may
replace $|I|$ by $|I\cap U|+(k-1-|A_1\cap U|)$, as in Lemma~\ref{lem:Ibound}. After this substitution the
quantity $F_0=h^\sharp-\mathrm{target}$ is a quadratic polynomial in the
$71$ variables with coefficients in $\mathbb Z[k]$.
\begin{enumerate}[label=(\arabic*)]
\item Of the $2{,}556$ possible quadratic monomials, $126$ have identically zero coefficient. For every integer $k\ge6$, the remaining $2{,}430$ coefficients split as follows: $2{,}348$ are nonnegative and $82$ are nonpositive. The two sets of monomials do not depend on $k$: this is checked at $k=6,7,8,9$, and
  on $[10,\infty)$ no coefficient changes sign. We drop the nonnegative terms. For the nonpositive ones we use $c\,xy\ge c\,k\,y$, which holds because
  $c\le0$ and $0\le x\le k$. This gives $F_0\ge c_0+\sum_v \ell'_v v$ with
  $c_0=-2(k-1)(2k^2-10k+13)$.
\item Every configuration satisfies four linear constraints $Mv\ge\1$:
  \begin{itemize}
  \item $b=\sum Z\ge1$;
  \item for each $X\in\{A_1,A_3,R\}$, $\#\{c\in U: c\notin X\}-b\ge1$, because $|U|=k+b$ and $|X|=k-1$.
  \end{itemize}
\item For $k\ge 9$, the multipliers
\[
\begin{aligned}
y_b&=\tfrac{1}{10}\bigl(9k^3+9k^2+1867k-17997\bigr), & y_{A_1}&=0,\\
y_{A_3}&=\tfrac{1}{10}\bigl(17k^3-133k^2-869k+10788\bigr), & y_R&=\tfrac{1}{10}\bigl(17k^3-116k^2-537k+6949\bigr)
\end{aligned}
\]
  are nonnegative and satisfy $\ell'_v\ge(M^{\mathsf T}y)_v$ for all $71$ variables, and
  $c_0+\sum y=k(3k^2+1)/10>0$. For $k=6,7,8$, the integer multipliers of Table~\ref{tab:fixedk} do
  the same.
\end{enumerate}
Consequently, since $y\ge0$ and $Mv\ge\1$,
\[
F_0\ \ge\ c_0+\ell'\!\cdot v\ \ge\ c_0+y\cdot Mv\ \ge\ c_0+\textstyle\sum y>0.
\]
For $k\ge10$ the $71$ inequalities and the sign conditions are identities in $\mathbb Q[k]$. They are
decided exactly by positivity of the leading coefficient together with real-root isolation on
$[10,\infty)$. For $k=6,\dots,9$ they are finitely many rational inequalities, evaluated exactly. See
Appendix~\ref{app:ues}.
\end{proof}

\subsection{Proof of Theorem~\ref{thm:house}}
By Lemma~\ref{lem:size} we may assume \eqref{eq:sizes}. For $k\ge6$ the three cases are covered by
Propositions~\ref{prop:equalpairs}, \ref{prop:es} and~\ref{prop:ues}. For $k\in\{4,5\}$ the whole
statement is verified by the complete type enumeration of Section~\ref{sec:computation}. (That
enumeration also re-proves the equal-pair and equal-shoulder cases for these $k$.) \qed

% =====================================================================
\section{Localization and compensation}\label{sec:pair}

\begin{lemma}[Localization]\label{lem:local}
Let $k\ge4$, and suppose $\delta_x>0$ for some $x\in A\cap B$. Then there are a $(k-1)$-set $C$, a colour
$\gamma\notin C\cup\{x\}$ and a colour $c_0\in C$ such that
\[
L(e_1)=L(e_3)=C\cup\{x\},\quad L(e_2)=L(e_4)=C\cup\{\gamma\},\quad B=\{x,\gamma\}\cup(C\setminus\{c_0\}),
\]
and $\delta_x=\sigma_k$.
\end{lemma}
\begin{proof}
By Lemma~\ref{lem:rowhouse}, $S(x)$ is a house count with roof $B\setminus\{x\}$. A positive deficit means
$S(x)<\mathrm{target}$. By Theorem~\ref{thm:house}, both base lists $L(e_i)\setminus\{x\}$ have size
$k-1$, so $x\in L(e_1)\cap L(e_3)$, and the house assignment is $\Delta$. We have $\gamma\ne x$ because
$\gamma\in B\setminus\{x\}$. The value of $\delta_x$ comes from Proposition~\ref{prop:equalpairs}.
\end{proof}

\begin{lemma}[Compensation]\label{lem:pair}
In the configuration of Lemma~\ref{lem:local}, every $x'\in B\setminus\{x\}$ satisfies
$\delta_x+\delta_{x'}<0$.
\end{lemma}
\begin{proof}
Put $P=C\cup\{x\}$ and $V=C\cup\{\gamma\}$, so that $L(e_1)=L(e_3)=P$ and $L(e_2)=L(e_4)=V$. By the
inclusion--exclusion count in the proof of Proposition~\ref{prop:equalpairs}, for colours $x'\ne y$
\[
F(x',y)=p(p-1)q(q-1)-2t(p-1)(q-1)+t(t-1),
\]
where $p=|P\setminus\{x'\}|$, $q=|V\setminus\{y\}|$ and $t=|(P\cap V)\setminus\{x',y\}|$. The other rows of $B$ are $\gamma$ and the $k-2$ colours
$c\in C\setminus\{c_0\}$. Summing the kernel over $y\in B\setminus\{x'\}$ gives
\[
\delta_\gamma=-2(k^4-5k^3+10k^2-8k+1),\qquad \delta_c=-2(2k^3-13k^2+31k-27).
\]
Hence
\[
\delta_x+\delta_\gamma=-2(k^4-5k^3+9k^2-3k-5)<0\quad\text{and}\quad
\delta_x+\delta_c=-2(2k^3-14k^2+36k-33)<0\quad(k\ge4).
\]
Both polynomials have positive coefficients after the substitution $k=4+m$.
\end{proof}

% =====================================================================
\section{Proof of Theorem~\ref{thm:main}}\label{sec:proof}

For $k=3$ use Corollary~\ref{cor:easy}(i). Let $k\ge4$ and put $q=|A\cap B|$. By
Corollary~\ref{cor:easy}(ii), rows $x\in A\setminus B$ have
$\delta_x\le-(k-1)(k-2)(k^2-7k+13)<0$. By Theorem~\ref{thm:house} and Lemma~\ref{lem:rowhouse}, rows
$x\in A\cap B$ have $\delta_x\le\sigma_k$. By Lemma~\ref{lem:local}, at most one row is positive, because
a positive row determines the configuration, and in it every other row of $B$ is negative by
Lemma~\ref{lem:pair}.
\begin{itemize}
\item If $q=0$, every term of \eqref{eq:deficits} is negative.
\item If $q=1$, the single inside row contributes at most $\sigma_k$. There are $k-1\ge1$ outside rows,
  each contributing at most $-(k-1)(k-2)(k^2-7k+13)$, and
  $(k-1)(k-2)(k^2-7k+13)-\sigma_k=(k-2)(k^3-8k^2+18k-7)>0$ for $k\ge4$.
\item If $q\ge2$ and some inside row $x$ is positive, pair it with any other inside row $x'$. Then
  $\delta_x+\delta_{x'}<0$ by Lemma~\ref{lem:pair}, and all remaining rows are nonpositive. If no row is
  positive, the sum is nonpositive trivially.
\end{itemize}
In every case $\sum_{x\in A}\delta_x\le0$, which is \eqref{eq:deficits}. \qed

% =====================================================================
\section{Computation and reproducibility}\label{sec:computation}

\subsection{\texorpdfstring{Complete type enumeration for $k\in\{4,5\}$}{Complete type enumeration for k=4,5}}
By \eqref{eq:shoulder}, the house count under \eqref{eq:sizes} depends only on
the following data:
\begin{itemize}
\item the number of colours of $U=L_2\cup L_4$ of each of the $24$ types, where a type records
  membership in $L_2$ only, $L_4$ only or both, and membership in each of $R$, $A_1$, $A_3$;
\item the overlap $|I\setminus U|$, which ranges over $0,\dots,\min\{|A_1\setminus U|,|A_3\setminus U|\}$.
\end{itemize}
Every feasible value enumerated is realized with fresh colours, so enumerating it covers every house
assignment in every colour universe. The proof needs this only for $k=4$ and $k=5$, where no multipliers of the kind used in
Proposition~\ref{prop:ues} exist (the corresponding dual linear program is infeasible). As a consistency check we ran it through
$k=9$.

\begin{center}\small
\begin{tabular}{@{}lrrrrrr@{}}
\toprule
$k$ & 4 & 5 & 6 & 7 & 8 & 9\\
\midrule
type configurations & 49{,}394 & 519{,}800 & 4{,}231{,}446 & 28{,}203{,}597 & 159{,}948{,}947 & 793{,}435{,}085\\
$\min(h-\mathrm{target})$ & $-4$ & $-12$ & $-24$ & $-40$ & $-60$ & $-84$\\
configurations with $h<\mathrm{target}$ & 1 & 1 & 1 & 1 & 1 & 1\\
\bottomrule
\end{tabular}
\end{center}

In each case $\min(h-\mathrm{target})=-\sigma_k$, and the unique configuration below the target is
$\Delta$. For $6\le k\le 9$ this agrees with the analytic proof. The proof itself needs the complete
enumeration only for $k=4,5$. The public reproducibility package supplies the complete enumerator and
recorded outputs for those two values. During development the same census was also run through $k=9$,
and two further independent implementations were used as cross-checks on overlapping ranges; all agreed
with the values in the table. Those additional audit implementations are not part of the public release.

\subsection{Independent verification}
The proof-critical computations were cross-checked independently during development. The public
reproducibility package contains a direct rebuild of the dual certificate, symbolic and small-case
checks of the displayed identities, the complete type enumeration required for $k=4,5$, and direct
small-universe checks of the full theorem. Additional independent audit implementations were used
during development but are not included in the public release. All checks agree:
\begin{itemize}
\item the counting formula \eqref{eq:shoulder}, compared with brute-force counts of house colourings;
\item Proposition~\ref{prop:es}, rebuilt symbolically, with the same table of coefficients and the same
  minimal patterns;
\item Proposition~\ref{prop:ues}. The polynomial certificate was re-derived twice, each time including
  the $2{,}348/82$ split of the quadratic coefficients and the exact dual inequalities. The certificates
  of Table~\ref{tab:fixedk} were checked by rebuilding the $71$ columns of $M$ from the definition in
  Appendix~\ref{app:ues};
\item the type enumeration of Section~\ref{sec:computation}: the supplied implementation reproduces
  the proof-critical $k=4,5$ cases, while the additional development implementations agreed on their
  overlapping ranges;
\item the identities in Lemma~\ref{lem:cyclebound}, Corollary~\ref{cor:easy},
  Proposition~\ref{prop:equalpairs} and Lemma~\ref{lem:pair}, re-derived symbolically and compared
  with brute-force counts for small $k$;
\item the whole of Theorem~\ref{thm:main}, compared with direct brute-force counts for normalized
  small colour universes at $k=3,4,5$; the supplied program and recorded outputs attain respectively
  $c_3=12$, $c_4=264$ and $c_5=1920$, with no smaller count.
\end{itemize}

\subsection*{Data and code availability}
The manuscript source, verification programs, recorded outputs and reproduction commands are publicly
available at \url{https://github.com/DMDTague/TriangularPrismChoosability}. The archival software
release will be tagged \texttt{v1.0.0} before submission. Its Zenodo DOI will be inserted here after
the release has been archived.

% =====================================================================
\section{Discussion}\label{sec:discussion}

The prism is $L(K_{2,3})=L(\Theta(2,2,2))$. Two natural families contain it:
\begin{itemize}
\item the line graphs $L(K_{2,n})=K_n\square K_2$, which bear on Question~\ref{q:KN} for the pairs
  $(K_n,K_2)$;
\item the line graphs of the theta graphs $\Theta(2,2,\ell)$ with $\ell$ even.
\end{itemize}
The gluing identity of Lemma~\ref{lem:glue} extends formally to both families. For
$\Theta(2,2,\ell)$ the path $uwv$ is replaced by a longer path, and for $K_{2,n}$ the $4$-cycle kernel is
replaced by the kernel of $K_{2,n-1}$. The difficulty lies in finding the analogue of the House Lemma. We leave both families to future work. More generally, one may ask whether
$L(R)$ is enumeratively chromatic-choosable for every bipartite graph $R$; at $m=\Delta(R)$ this is the
per-edge question of~\cite{TagueExact}. By Galvin's theorem this
would be an enumerative strengthening of the list edge-colouring theorem for bipartite graphs.
The general thresholds \cite{DZ23,ZD26,KirovLean} cover large list sizes. We found no result that covers
the critical case $m=\Delta(R)$ in general.

\appendix
\section{The certificate of Proposition~\ref{prop:ues}}\label{app:ues}

\paragraph{Multiplier provenance.}
For fixed $k$, the four dual multipliers solve a rational linear feasibility problem with the
$71$ inequalities $M^\top y\le\ell'$. The polynomial certificate was obtained from exact rational
solutions at several integer values in the stable range and interpolation in $k$. This discovery
step is not part of the proof: for every $k\ge10$, nonnegativity of the multipliers and all $71$
residual inequalities are verified independently by exact polynomial arithmetic and real-root
isolation on $[10,\infty)$.

The constraint matrix $M$ has one row for each constraint $b\ge1$, $A_1$-miss, $A_3$-miss and $R$-miss.
Encode a type $\tau$ as $(\tau_R,\tau_{A_1},\tau_{A_3})$. The column of $D_\tau$ is
$(0,\,1-\tau_{A_1},\,1-\tau_{A_3},\,1-\tau_R)$. The column of $Z_{\tau,\tau'}$ is
\[
\bigl(1,\ (1-\tau_{A_1})+(1-\tau'_{A_1})-1,\ (1-\tau_{A_3})+(1-\tau'_{A_3})-1,\ (1-\tau_R)+(1-\tau'_R)-1\bigr).
\]
\begin{table}[ht]\centering\small
\begin{tabular}{@{}lrrrrrrl@{}}
\toprule
$k$ & $y_b$ & $y_{A_1}$ & $y_{A_3}$ & $y_R$ & $c_0$ & $c_0+\sum y$ & residuals equal to $0$\\
\midrule
6 & 66 & 0 & 87 & 130 & $-250$ & 33 & $D_{011},D_{100},Z_{011,101}$\\
7 & 218 & 13 & 175 & 252 & $-492$ & 166 & $D_{011},D_{100},Z_{000,000},Z_{000,001}$\\
8 & 519 & 8 & 316 & 288 & $-854$ & 277 & seven $Z$-variables$^{*}$\\
9 & $\tfrac{3048}{5}$ & 0 & $\tfrac{4587}{10}$ & $\tfrac{5113}{10}$ & $-1360$ & $\tfrac{1098}{5}$ & none (minimum $\tfrac{1113}{10}$)\\
\bottomrule
\end{tabular}
\caption{Certificates of Proposition~\ref{prop:ues} for $6\le k\le 9$. The row $k=9$ is the polynomial
certificate evaluated at $k=9$. Every residual $r_v=\ell'_v-(M^{\mathsf T}y)_v$ is nonnegative.
$^{*}$These are $Z_{000,000}$, $Z_{000,100}$, $Z_{010,000}$, $Z_{010,100}$, $Z_{011,101}$, $Z_{100,000}$ and $Z_{100,100}$.}
\label{tab:fixedk}
\end{table}

Order the variables as follows: first the $D_\tau$, then the $Z_{\tau,\tau'}$, each in lexicographic
order of the type strings $\tau=\tau_R\tau_{A_1}\tau_{A_3}$. A nonpositive quadratic term
$c\,v_iv_j$ with $i\le j$ is bounded below by $c\,k\,v_j$, and so it is charged to the later variable. The
coefficient $\ell'_v$ is the linear coefficient of $v$ in $F_0$ plus all charges made to $v$. The residual
is $r_v:=\ell'_v-(M^{\mathsf T}y)_v$.

The following table lists $\ell'_v$, $r_v$ and $r_v(10)$ for all $71$ variables. Each $r_v$ is a cubic
with positive leading coefficient and no real root in $[10,\infty)$; the largest real root of any
residual is about $8.35$. The minimum of $r_v(10)$ is $443/2$. Finally
$c_0+\sum_i y_i=k(3k^2+1)/10$. Every row of the table is checked by the script
\texttt{certificate\_rebuild.py} in the supplementary material.

% --- begin inlined ues_residuals.tex ---
{\footnotesize\begin{longtable}{@{}llll@{}}
\caption{Charged linear coefficients $\ell'_v$ and dual residuals $r_v=\ell'_v-(M^{\mathsf T}y)_v$ for Proposition~\ref{prop:ues}.}\\
\toprule $v$ & $\ell'_v(k)$ & $10\,r_v(k)$ & $r_v(10)$\\ \midrule \endhead
$D_{000}$ & $4 k^{3} - 24 k^{2} + 43 k - 19$ & $6 k^{3} + 9 k^{2} + 1836 k - 17927$ & $7333 / 10$\\
$D_{001}$ & $3 k^{3} - 19 k^{2} + 37 k - 18$ & $13 k^{3} - 74 k^{2} + 907 k - 7129$ & $7541 / 10$\\
$D_{010}$ & $4 k^{3} - 23 k^{2} + 40 k - 18$ & $6 k^{3} + 19 k^{2} + 1806 k - 17917$ & $8043 / 10$\\
$D_{011}$ & $2 k^{3} - 12 k^{2} + 24 k - 14$ & $3 k^{3} - 4 k^{2} + 777 k - 7089$ & $3281 / 10$\\
$D_{100}$ & $2 k^{3} - 14 k^{2} + 29 k - 15$ & $3 k^{3} - 7 k^{2} + 1159 k - 10938$ & $1476 / 5$\\
$D_{101}$ & $k^{3} - 8 k^{2} + 19 k - 12$ & $10 k^{3} - 80 k^{2} + 190 k - 120$ & $378$\\
$D_{110}$ & $2 k^{3} - 12 k^{2} + 22 k - 12$ & $3 k^{3} + 13 k^{2} + 1089 k - 10908$ & $2141 / 5$\\
$Z_{000,000}$ & $5 k^{3} - 28 k^{2} + 48 k - 21$ & $7 k^{3} - 40 k^{2} + 19 k + 50$ & $324$\\
$Z_{000,001}$ & $4 k^{3} - 24 k^{2} + 44 k - 21$ & $14 k^{3} - 133 k^{2} - 890 k + 10838$ & $1319 / 5$\\
$Z_{000,010}$ & $6 k^{3} - 32 k^{2} + 51 k - 21$ & $17 k^{3} - 80 k^{2} + 49 k + 50$ & $954$\\
$Z_{000,011}$ & $4 k^{3} - 23 k^{2} + 42 k - 21$ & $14 k^{3} - 123 k^{2} - 910 k + 10838$ & $1719 / 5$\\
$Z_{000,100}$ & $4 k^{3} - 24 k^{2} + 44 k - 21$ & $14 k^{3} - 116 k^{2} - 558 k + 6999$ & $3819 / 10$\\
$Z_{000,101}$ & $3 k^{3} - 19 k^{2} + 36 k - 21$ & $21 k^{3} - 199 k^{2} - 1507 k + 17787$ & $3817 / 10$\\
$Z_{000,110}$ & $5 k^{3} - 28 k^{2} + 47 k - 21$ & $24 k^{3} - 156 k^{2} - 528 k + 6999$ & $10119 / 10$\\
$Z_{000,111}$ & $3 k^{3} - 18 k^{2} + 33 k - 21$ & $21 k^{3} - 189 k^{2} - 1537 k + 17787$ & $4517 / 10$\\
$Z_{001,000}$ & $5 k^{3} - 28 k^{2} + 48 k - 21$ & $24 k^{3} - 173 k^{2} - 850 k + 10838$ & $4519 / 5$\\
$Z_{001,001}$ & $4 k^{3} - 24 k^{2} + 44 k - 21$ & $31 k^{3} - 266 k^{2} - 1759 k + 21626$ & $4218 / 5$\\
$Z_{001,010}$ & $6 k^{3} - 32 k^{2} + 51 k - 21$ & $34 k^{3} - 213 k^{2} - 820 k + 10838$ & $7669 / 5$\\
$Z_{001,011}$ & $4 k^{3} - 23 k^{2} + 42 k - 21$ & $31 k^{3} - 256 k^{2} - 1779 k + 21626$ & $4618 / 5$\\
$Z_{001,100}$ & $4 k^{3} - 24 k^{2} + 44 k - 21$ & $31 k^{3} - 249 k^{2} - 1427 k + 17787$ & $9617 / 10$\\
$Z_{001,101}$ & $3 k^{3} - 19 k^{2} + 33 k - 21$ & $38 k^{3} - 332 k^{2} - 2406 k + 28575$ & $1863 / 2$\\
$Z_{001,110}$ & $5 k^{3} - 28 k^{2} + 47 k - 21$ & $41 k^{3} - 289 k^{2} - 1397 k + 17787$ & $15917 / 10$\\
$Z_{001,111}$ & $3 k^{3} - 18 k^{2} + 32 k - 21$ & $38 k^{3} - 322 k^{2} - 2416 k + 28575$ & $2043 / 2$\\
$Z_{010,000}$ & $5 k^{3} - 28 k^{2} + 47 k - 21$ & $7 k^{3} - 40 k^{2} + 9 k + 50$ & $314$\\
$Z_{010,001}$ & $4 k^{3} - 24 k^{2} + 43 k - 21$ & $14 k^{3} - 133 k^{2} - 900 k + 10838$ & $1269 / 5$\\
$Z_{010,010}$ & $6 k^{3} - 32 k^{2} + 50 k - 21$ & $17 k^{3} - 80 k^{2} + 39 k + 50$ & $944$\\
$Z_{010,011}$ & $4 k^{3} - 23 k^{2} + 41 k - 21$ & $14 k^{3} - 123 k^{2} - 920 k + 10838$ & $1669 / 5$\\
$Z_{010,100}$ & $4 k^{3} - 24 k^{2} + 43 k - 21$ & $14 k^{3} - 116 k^{2} - 568 k + 6999$ & $3719 / 10$\\
$Z_{010,101}$ & $3 k^{3} - 19 k^{2} + 36 k - 21$ & $21 k^{3} - 199 k^{2} - 1507 k + 17787$ & $3817 / 10$\\
$Z_{010,110}$ & $5 k^{3} - 28 k^{2} + 46 k - 21$ & $24 k^{3} - 156 k^{2} - 538 k + 6999$ & $10019 / 10$\\
$Z_{010,111}$ & $3 k^{3} - 18 k^{2} + 33 k - 21$ & $21 k^{3} - 189 k^{2} - 1537 k + 17787$ & $4517 / 10$\\
$Z_{011,000}$ & $4 k^{3} - 23 k^{2} + 42 k - 21$ & $14 k^{3} - 123 k^{2} - 910 k + 10838$ & $1719 / 5$\\
$Z_{011,001}$ & $3 k^{3} - 19 k^{2} + 36 k - 21$ & $21 k^{3} - 216 k^{2} - 1839 k + 21626$ & $1318 / 5$\\
$Z_{011,010}$ & $5 k^{3} - 27 k^{2} + 45 k - 21$ & $24 k^{3} - 163 k^{2} - 880 k + 10838$ & $4869 / 5$\\
$Z_{011,011}$ & $3 k^{3} - 18 k^{2} + 35 k - 21$ & $21 k^{3} - 206 k^{2} - 1849 k + 21626$ & $1768 / 5$\\
$Z_{011,100}$ & $3 k^{3} - 19 k^{2} + 38 k - 21$ & $21 k^{3} - 199 k^{2} - 1487 k + 17787$ & $4017 / 10$\\
$Z_{011,101}$ & $2 k^{3} - 14 k^{2} + 12 k - 21$ & $28 k^{3} - 282 k^{2} - 2616 k + 28575$ & $443 / 2$\\
$Z_{011,110}$ & $4 k^{3} - 23 k^{2} + 41 k - 21$ & $31 k^{3} - 239 k^{2} - 1457 k + 17787$ & $10317 / 10$\\
$Z_{011,111}$ & $2 k^{3} - 13 k^{2} + 21 k - 21$ & $28 k^{3} - 272 k^{2} - 2526 k + 28575$ & $823 / 2$\\
$Z_{100,000}$ & $4 k^{3} - 24 k^{2} + 44 k - 21$ & $14 k^{3} - 116 k^{2} - 558 k + 6999$ & $3819 / 10$\\
$Z_{100,001}$ & $3 k^{3} - 20 k^{2} + 38 k - 21$ & $21 k^{3} - 209 k^{2} - 1487 k + 17787$ & $3017 / 10$\\
$Z_{100,010}$ & $5 k^{3} - 28 k^{2} + 47 k - 21$ & $24 k^{3} - 156 k^{2} - 528 k + 6999$ & $10119 / 10$\\
$Z_{100,011}$ & $3 k^{3} - 19 k^{2} + 38 k - 21$ & $21 k^{3} - 199 k^{2} - 1487 k + 17787$ & $4017 / 10$\\
$Z_{100,100}$ & $3 k^{3} - 20 k^{2} + 40 k - 21$ & $21 k^{3} - 192 k^{2} - 1135 k + 13948$ & $2199 / 5$\\
$Z_{100,101}$ & $2 k^{3} - 15 k^{2} + 28 k - 21$ & $28 k^{3} - 275 k^{2} - 2124 k + 24736$ & $1998 / 5$\\
$Z_{100,110}$ & $4 k^{3} - 24 k^{2} + 43 k - 21$ & $31 k^{3} - 232 k^{2} - 1105 k + 13948$ & $5349 / 5$\\
$Z_{100,111}$ & $2 k^{3} - 14 k^{2} + 31 k - 21$ & $28 k^{3} - 265 k^{2} - 2094 k + 24736$ & $2648 / 5$\\
$Z_{101,000}$ & $4 k^{3} - 24 k^{2} + 44 k - 21$ & $31 k^{3} - 249 k^{2} - 1427 k + 17787$ & $9617 / 10$\\
$Z_{101,001}$ & $3 k^{3} - 20 k^{2} + 38 k - 21$ & $38 k^{3} - 342 k^{2} - 2356 k + 28575$ & $1763 / 2$\\
$Z_{101,010}$ & $5 k^{3} - 28 k^{2} + 47 k - 21$ & $41 k^{3} - 289 k^{2} - 1397 k + 17787$ & $15917 / 10$\\
$Z_{101,011}$ & $3 k^{3} - 19 k^{2} + 38 k - 21$ & $38 k^{3} - 332 k^{2} - 2356 k + 28575$ & $1963 / 2$\\
$Z_{101,100}$ & $3 k^{3} - 20 k^{2} + 40 k - 21$ & $38 k^{3} - 325 k^{2} - 2004 k + 24736$ & $5098 / 5$\\
$Z_{101,101}$ & $2 k^{3} - 15 k^{2} + 28 k - 21$ & $45 k^{3} - 408 k^{2} - 2993 k + 35524$ & $4897 / 5$\\
$Z_{101,110}$ & $4 k^{3} - 24 k^{2} + 43 k - 21$ & $48 k^{3} - 365 k^{2} - 1974 k + 24736$ & $8248 / 5$\\
$Z_{101,111}$ & $2 k^{3} - 14 k^{2} + 31 k - 21$ & $45 k^{3} - 398 k^{2} - 2963 k + 35524$ & $5547 / 5$\\
$Z_{110,000}$ & $4 k^{3} - 23 k^{2} + 41 k - 21$ & $14 k^{3} - 106 k^{2} - 588 k + 6999$ & $4519 / 10$\\
$Z_{110,001}$ & $3 k^{3} - 19 k^{2} + 37 k - 21$ & $21 k^{3} - 199 k^{2} - 1497 k + 17787$ & $3917 / 10$\\
$Z_{110,010}$ & $5 k^{3} - 27 k^{2} + 44 k - 21$ & $24 k^{3} - 146 k^{2} - 558 k + 6999$ & $10819 / 10$\\
$Z_{110,011}$ & $3 k^{3} - 18 k^{2} + 35 k - 21$ & $21 k^{3} - 189 k^{2} - 1517 k + 17787$ & $4717 / 10$\\
$Z_{110,100}$ & $3 k^{3} - 19 k^{2} + 37 k - 21$ & $21 k^{3} - 182 k^{2} - 1165 k + 13948$ & $2549 / 5$\\
$Z_{110,101}$ & $2 k^{3} - 14 k^{2} + 31 k - 21$ & $28 k^{3} - 265 k^{2} - 2094 k + 24736$ & $2648 / 5$\\
$Z_{110,110}$ & $4 k^{3} - 23 k^{2} + 40 k - 21$ & $31 k^{3} - 222 k^{2} - 1135 k + 13948$ & $5699 / 5$\\
$Z_{110,111}$ & $2 k^{3} - 13 k^{2} + 28 k - 21$ & $28 k^{3} - 255 k^{2} - 2124 k + 24736$ & $2998 / 5$\\
$Z_{111,000}$ & $3 k^{3} - 18 k^{2} + 35 k - 21$ & $21 k^{3} - 189 k^{2} - 1517 k + 17787$ & $4717 / 10$\\
$Z_{111,001}$ & $2 k^{3} - 14 k^{2} + 24 k - 21$ & $28 k^{3} - 282 k^{2} - 2496 k + 28575$ & $683 / 2$\\
$Z_{111,010}$ & $4 k^{3} - 22 k^{2} + 38 k - 21$ & $31 k^{3} - 229 k^{2} - 1487 k + 17787$ & $11017 / 10$\\
$Z_{111,011}$ & $2 k^{3} - 13 k^{2} + 29 k - 21$ & $28 k^{3} - 272 k^{2} - 2446 k + 28575$ & $983 / 2$\\
$Z_{111,100}$ & $2 k^{3} - 14 k^{2} + 31 k - 21$ & $28 k^{3} - 265 k^{2} - 2094 k + 24736$ & $2648 / 5$\\
$Z_{111,101}$ & $k^{3} - 9 k^{2} + 4 k - 21$ & $35 k^{3} - 348 k^{2} - 3233 k + 35524$ & $1697 / 5$\\
$Z_{111,110}$ & $3 k^{3} - 18 k^{2} + 34 k - 21$ & $38 k^{3} - 305 k^{2} - 2064 k + 24736$ & $5798 / 5$\\
$Z_{111,111}$ & $k^{3} - 8 k^{2} + 22 k - 21$ & $35 k^{3} - 338 k^{2} - 3053 k + 35524$ & $3097 / 5$\\
\bottomrule
\end{longtable}}
% --- end inlined ues_residuals.tex ---

\begin{thebibliography}{99}

\bibitem{AM} S.~Allred and J.~A.~Mudrock, \emph{Enumerative chromatic choosability}, preprint,
  arXiv:2505.05662, 2025.

\bibitem{CLMNW} Y.~Chi, S.~Lee, F.~Morrissette, J.~A.~Mudrock, G.~Nguyen and B.~Whatley,
  \emph{Enumeratively chromatic-choosable theta graphs}, Enumer.\ Combin.\ Appl.\ \textbf{6}(3) (2026),
  Article S3R23; arXiv:2605.10861.

\bibitem{DZ23} F.~Dong and M.~Zhang, \emph{An improved lower bound of $P(G,L)-P(G,k)$ for
  $k$-assignments $L$}, J.\ Combin.\ Theory Ser.\ B \textbf{161} (2023), 109--119; arXiv:2206.14536.

\bibitem{Donner} Q.~Donner, \emph{On the number of list-colorings}, J.\ Graph Theory \textbf{16}
  (1992), 239--245.

\bibitem{Galvin} F.~Galvin, \emph{The list chromatic index of a bipartite multigraph}, J.\ Combin.\
  Theory Ser.\ B \textbf{63} (1995), 153--158.

\bibitem{KN} R.~Kirov and R.~Naimi, \emph{List coloring and $n$-monophilic graphs}, Ars Combin.\
  \textbf{124} (2016), 329--340; arXiv:1004.5183.

\bibitem{KirovLean} R.~Kirov, \emph{enumerative-chromatic-choosability}: a Lean~4 formalization of the
  list colour function, \url{https://github.com/rkirov/enumerative-chromatic-choosability}, accessed
  27~September 2026 (commit \texttt{95d4009}).

\bibitem{KS} A.~V.~Kostochka and A.~F.~Sidorenko, \emph{Problem session}, in: Fourth Czechoslovak
  Symposium on Combinatorics, Graphs and Complexity (J.~Ne\v{s}et\v{r}il and M.~Fiedler, eds.), Ann.\
  Discrete Math.\ \textbf{51}, North-Holland, Amsterdam, 1992, p.~380.

\bibitem{TagueExact} D.~Tague, \emph{Exact list-edge labelings of bipartite multigraphs: uncrossing,
  nested minimizers, and a $2$-factor criterion for equality}, preprint, 2026.

\bibitem{Thomassen} C.~Thomassen, \emph{The chromatic polynomial and list colorings}, J.\ Combin.\
  Theory Ser.\ B \textbf{99} (2009), 474--479.

\bibitem{WQY} W.~Wang, J.~Qian and Z.~Yan, \emph{When does the list-coloring function of a graph equal
  its chromatic polynomial?}, J.\ Combin.\ Theory Ser.\ B \textbf{122} (2017), 543--549.

\bibitem{ZD26} M.~Zhang and F.~Dong, \emph{When chromatic polynomials coincide with list-color
  functions: a threshold linear in the maximum degree}, preprint, arXiv:2609.08540, 2026.

\bibitem{ZD26np} M.~Zhang and F.~Dong, \emph{Non-persistence of equality between chromatic polynomials
  and list-color functions}, preprint, arXiv:2608.19773, 2026.

\end{thebibliography}
\end{document}